\documentclass[12pt,a4paper,UTF8]{ctexart}

% 必要宏包
\usepackage{geometry}
\usepackage{amsmath}
\usepackage{amssymb}
\usepackage{graphicx}
\usepackage{booktabs}
\usepackage{hyperref}
\usepackage{float}
\usepackage{enumitem}
\usepackage{multirow}
\usepackage{array}

% 页面布局
\geometry{left=3cm,right=2.5cm,top=2.5cm,bottom=2.5cm}

% 超链接样式
\hypersetup{
    colorlinks=true,
    linkcolor=blue,
    citecolor=blue,
    urlcolor=blue
}

% 论文信息
\title{\textbf{基于耦合多智能体强化学习的\\低空航空交通协同管控模型}}
\author{作者：彭义森 \quad 指导老师：韩云祥 \\ 四川大学}
\date{2026年7月}

\begin{document}

\maketitle

\begin{abstract}
随着先进空中交通（AAM）的快速发展，电动垂直起降飞行器（eVTOL）在城市低空的高密度运行对传统空中交通管理提出了严峻挑战。现有多智能体强化学习（MARL）方法在时空耦合建模、稀疏奖励处理与物理约束建模三个维度上均存在不足。针对城市低空十字交叉路口的多飞行器冲突避免与流量协同优化问题，本文提出一种基于耦合MARL的低空交通管控模型。模型核心包括：动态图注意力网络（DyGAT）融合空间注意力、时序注意力与基于物理量的动态边权重，刻画飞行器间时空耦合关联；分层奖励函数兼顾飞行安全与运行效率；课程学习与基于冲突事件的轨迹加权策略解决稀疏奖励问题。在融合风扰与传感器噪声的仿真环境中，所提模型在21架eVTOL交叉路口场景中取得了最优性能：相比RR方法，冲突次数降低57.2\%，平均运行延迟缩短63.9\%；相比当前最优的CTDE-MARL方法MAPPO，冲突次数降低10.3\%（$p=0.180$，未达统计显著），无冲突完成率从57.8\%提升至68.4\%（$p<0.001$），平均运行延迟缩短12.3\%。当前仿真环境为二维单交叉路口的简化抽象，与真实三维广域低空交通存在差距；实验尚未包含MPC、ORCA与CBF等工程安全关键方法的对比。本文工作为城市低空交通的分布式自主管控提供了一条有潜力的技术路线。

\noindent \textbf{关键词：} 低空航空交通管理；多智能体强化学习；动态图注意力网络；冲突避免；集中训练分布式执行；时空耦合建模
\end{abstract}

%==========================================================================
\section{引言}
%==========================================================================

城市地面交通拥堵的日益加剧推动了以eVTOL为核心的城市空中交通（UAM）产业快速发展。FAA预测，到2040年将有数百万架无人机与eVTOL在美国低空空域运行。与传统民航航路的结构化交通不同，城市低空交通具有高密度、强动态、高异构性特征，交叉路口与汇合点的协同管控与冲突避免已成为低空无人机交通管理系统（UTM）面临的核心挑战。

传统空中交通流量管理（ATFM）多采用集中式优化架构（整数规划、蒙特卡洛仿真等），依赖预先规划的飞行计划与静态航路设计，难以适配低空环境中的动态不确定性。近年来，MARL因其分布式决策与自适应学习特性被成功引入交通管控领域。Agogino与Tumer的系列工作率先为ATFM建立了多智能体学习范式：从分布式航路点智能体调控框架\cite{ref1}，到基于差异奖励的大规模流量管理\cite{ref2}，系统性地验证了MARL的有效性。Brittain与Wei设计了基于A2C结合PPO损失函数的深度MARL框架，在BlueSky仿真器中实现了99.97\%的冲突规避率\cite{ref3}。Ghosh等提出深度集成MARL方法，通过元策略仲裁不同专家策略\cite{ref9}。Xie等将强化学习与遗传算法结合，构建了面向UAM需求-容量平衡的混合学习框架\cite{ref6}。Li等最近提出的MS-GRL采用多尺度图增强强化学习解决密集无人机网络中的冲突解决问题，在100架UAV规模下验证了图结构对冲突规避的关键作用\cite{ref24}。

尽管上述工作取得了显著进展，现有研究在三个维度上仍存在不足：

\textbf{时空耦合建模方面：}多数方法通过全连接网络或静态图结构编码智能体关联。Brittain与Wei仅考虑最近邻静态关联\cite{ref3}；Spatharis等通过预定义关系矩阵编码交互\cite{ref7}；Ghosh等依赖网格化Geo-hash密度图\cite{ref9}；MS-GRL虽采用了多尺度图注意力机制，但其图结构基于固定阈值而非物理量驱动的连续动态权重\cite{ref24}。这些方法的共同局限在于：当飞行器快速移动导致邻域拓扑剧烈变化时，静态或半静态图结构无法及时反映飞行器间动态变化的冲突风险等级。

\textbf{训练效率方面：}低空交通中冲突事件属于低概率高风险事件，训练过程中冲突样本占比通常低于0.1\%。Aslani等指出标准$\epsilon$-greedy探索策略在高维连续动作空间中难以高效覆盖稀疏关键状态区域\cite{ref10}；翟子洋等的综述将稀疏奖励列为MARL交通信号控制的首要挑战\cite{ref11}。标准经验回放机制中冲突样本采样概率极低，且冲突规避的奖励信号仅在冲突发生或成功通过时才出现，导致信用分配困难。

\textbf{环境建模方面：}多数研究采用简化模型，忽略低空环境的复杂性与不确定性。Spatharis等指出传统固定安全距离无法适配不同飞行速度下的制动需求\cite{ref7}。大多数现有研究将飞行器视为质点，忽略最大加速度、最大转弯角等物理运动约束。此外，低空风场建模多采用均匀随机扰动，与真实Dryden风湍流谱或城市峡谷非均匀风场存在显著简化。需要指出，本文当前的风扰模型同样采用均匀随机分布——这一简化对结果的影响在第5.4节中诚实讨论。

针对上述问题，本文设计了一种耦合多智能体强化学习管控模型，核心思路是：(1)~通过DyGAT的物理量驱动动态边权重精准刻画飞行器间时空耦合关系；(2)~通过课程学习与冲突轨迹加权在保持PPO on-policy理论框架的前提下高效利用稀疏冲突样本；(3)~通过动态安全距离与物理约束建模提升环境的工程合理性。

本文主要贡献包括：(1)~构建了面向低空交叉路口冲突规避的仿真环境\texttt{AirTrafficEnv}，引入与速度相关的动态安全距离模型与eVTOL物理运动约束（参数取值为参考当前主流eVTOL机型公开技术指标的代表性设定，详见第3.2.2节）；(2)~提出基于DyGAT的耦合智能体群网络，通过距离、速度、航向角三个物理量直接计算动态边权重——消融实验表明该模块使冲突规避能力提升174\%；(3)~设计课程学习与基于冲突事件的episode级轨迹加权训练框架（$\lambda=2.0$），在不引入off-policy数据的前提下提升冲突样本利用效率；(4)~与8种基线方法进行系统对比，引入无冲突完成率（CFCR）等安全攸关指标，并通过统计显著性检验诚实呈现TC不显著（$p=0.180$）与CFCR高度显著（$p<0.001$）的差异图景。

%==========================================================================
\section{相关工作}
%==========================================================================

本节从ATFM方法与MARL在交通管控中的应用两个维度梳理相关研究，表\ref{tab:related_work}给出系统性对比。

\subsection{传统空中交通流量管理方法}

传统ATFM方法可分为集中式优化与分布式启发式两类。集中式优化方面，Bertsimas等采用整数规划求解空域流量分配，Menon等通过交通流欧拉模型的线性规划实现动态调控\cite{ref6}。此类方法在可预测场景中效果良好，但在高密度、强动态场景中计算复杂度高。启发式方法方面，Agogino与Tumer设计了基于规则的反馈控制机制\cite{ref4}，计算效率高但泛化能力受限。此外，MPC与速度障碍法（ORCA/RVO）\cite{ref22}在航空冲突避免领域已有广泛应用，CBF\cite{ref23}作为安全关键系统的形式化保证方法近年来受到广泛关注——受限于当前仿真环境的二维简化设定，本文尚未将上述方法纳入实验对比（详见第5.2节）。

\subsection{强化学习在交通管控中的应用}

在单智能体RL方面，DQN已被应用于单飞行器冲突规避与单交叉口信号控制等场景。Aslani等在城市地面交通信号控制中系统对比了DQN与A2C的鲁棒性\cite{ref10}。国内研究者在该方向也开展了大量工作：王迎\cite{ref13}、赖建辉\cite{ref16}、张新武\cite{ref17}分别在硕士论文中将DRL应用于交通信号控制优化；赵晓华等\cite{ref14}、承向军等\cite{ref15}、卢守峰等\cite{ref18}探索了Q-learning及其变体在交叉口信号控制中的应用；江波等将DQN应用于机场道口协同智能调速\cite{ref19}；张春阳与席雅琪综述了AI在城市交通控制中的应用现状\cite{ref8}。需要说明，地面交通信号控制中的MARL方法在稀疏奖励、信用分配与CTDE架构等方面的设计思路对低空交通管控具有方法论参考价值。

在多智能体RL领域，Ghavamzadeh等提出层级式MARL框架\cite{ref12}；Tumer与Agogino建立了基于航路点智能体的分布式ATM学习框架\cite{ref2,ref4,ref5}；Ghosh等提出深度集成MARL方法\cite{ref9}；翟子洋等的综述系统梳理了从独立Q-learning到QMIX、MAPPO等CTDE方法的演进路线\cite{ref11}。Li等的MS-GRL采用多尺度图注意力与Double DQN解决密集UAV网络冲突\cite{ref24}，与本文目标相近，但在方法上存在本质差异：MS-GRL使用Double DQN处理离散动作空间，图结构依赖固定距离阈值；本文采用PPO处理连续加速度控制，图结构由物理量驱动的连续动态权重定义。

\subsection{本文方法与现有工作的差异化定位}

表\ref{tab:related_work}给出系统性对比。现有工作在时空耦合建模方面多采用全连接或静态图结构，在训练机制方面未针对冲突样本稀疏进行专门设计，在环境建模方面多采用简化模型。本文方法在这些维度上均做出了针对性改进，但需要指出：表中所列的``完整''物理约束与``风扰+噪声''扰动建模仍是对真实低空环境的简化——Dryden标准风谱模型、城市峡谷效应等复杂气象因素尚未纳入。

\begin{table}[H]
    \centering
    \caption{本文方法与代表性相关工作的系统性对比}
    \label{tab:related_work}
    \resizebox{0.9\textwidth}{!}{%
    \begin{tabular}{|c|c|c|c|c|c|}
        \hline
        \textbf{方法} & \textbf{时空耦合建模} & \textbf{图结构类型} & \textbf{稀疏奖励处理} & \textbf{物理约束} & \textbf{扰动建模} \\
        \hline
        Brittain \& Wei \cite{ref3} & LSTM+FC & 静态全连接 & $\times$ & 部分 & $\times$ \\
        \hline
        Ghosh et al. \cite{ref9} & Geo-hash密度图 & 网格化 & $\times$ & 部分 & $\times$ \\
        \hline
        Spatharis et al. \cite{ref7} & 预定义关系矩阵 & 层级分组 & $\times$ & $\times$ & $\times$ \\
        \hline
        MS-GRL \cite{ref24} & 多尺度GAT & 固定阈值图 & $\times$ & 部分 & $\times$ \\
        \hline
        QMIX/MAPPO基线 & FC/RNN & 全连接 & $\times$ & 部分 & $\times$ \\
        \hline
        \textbf{本文方法} & \textbf{DyGAT(空间+时序+动态边)} & \textbf{动态图} & \textbf{课程学习+轨迹加权} & \textbf{完整} & \textbf{风扰+噪声} \\
        \hline
    \end{tabular}%
    }
\end{table}

%==========================================================================
\section{问题建模与管控模型设计}
%==========================================================================

本文聚焦城市低空十字交叉路口场景：$N$架eVTOL从不同方向飞向交叉中心，需在保持安全间隔的前提下以最小延迟通过交叉区域。本节将问题建模为Dec-MDP，依次阐述环境模型、奖励函数、DyGAT网络、单智能体网络、GAT-Mixer混合网络与训练优化机制。

\subsection{分布式马尔可夫决策过程建模}

将$N$架eVTOL的协同管控形式化定义为Dec-MDP六元组$\langle \mathcal{N}, \mathcal{S}, \mathcal{A}, \mathcal{P}, \mathcal{R}, \gamma \rangle$：

\begin{itemize}
    \item \textbf{智能体集合} $\mathcal{N} = \{1, 2, \ldots, N\}$：每架eVTOL为一个独立智能体，$N$由课程学习机制动态调整（5$\to$21）。
    \item \textbf{全局状态} $\mathcal{S}$：$S = (s_1, s_2, \ldots, s_N)$，$s_i$为飞行器$i$的原始状态$[x_i, y_i, v_i, v_i^{\text{target}}, d_i^{\text{center}}]$（5维）。该5维向量为Dec-MDP的形式化状态变量；智能体网络的实际输入是经邻域扩展与历史序列编码后的局部观测$o_i$（见第3.2.1节）。全局状态$S$仅在CTDE训练阶段被GAT-Mixer访问。
    \item \textbf{联合动作} $\mathcal{A}$：$A = (a_1, a_2, \ldots, a_N)$，$a_i \in \mathbb{R}$为连续加速度指令。
    \item \textbf{状态转移} $\mathcal{P}(S' \mid S, A)$：受飞行器运动学、最大转弯角约束、风扰与传感器噪声的共同影响。
    \item \textbf{联合奖励} $\mathcal{R}$：$R = \frac{1}{N}\sum_{i=1}^{N} r_i$，$r_i$由安全、速度、协同与效率四个分量构成（第3.3节）。
    \item \textbf{折扣因子} $\gamma = 0.99$。
\end{itemize}

优化目标为$\pi^* = \arg\max_\pi \mathbb{E}_{\pi}[\sum_{t=0}^{T} \gamma^t R_t]$。采用CTDE架构：训练阶段GAT-Mixer（第3.6节）利用全局状态学习协同策略；执行阶段各智能体仅基于局部观测独立决策。

\subsection{环境模型}

仿真环境\texttt{AirTrafficEnv}以二维平面模拟交叉路口空域（$1000\text{m} \times 1000\text{m}$），中心为坐标原点。飞行器从均匀分布于圆周的初始位置出发，方向指向中心。该环境为低空交叉路口的简化抽象——三维空间、多路口网络、标准风谱模型等维度的扩展是未来工作的重要方向。

\subsubsection{状态空间的三层定义}

状态信息区分为三个层次：\textbf{原始状态}（5维向量$[x_i, y_i, v_i, v_i^{\text{target}}, d_i^{\text{center}}]$，即Dec-MDP的形式化状态$s_i$）；\textbf{局部观测}（在原始状态基础上经500m邻域半径与$T=10$时间窗口扩展，包含邻域飞行器的相对位置、相对速度、加速度、航向角及历史轨迹序列）；\textbf{图节点特征}（局部观测经LSTM编码与DyGAT聚合后得到的256维时空耦合特征向量$h_i^{\text{coupled}}$）。端到端流程为：$s_i^{\text{raw}} \to o_i \to$ LSTM编码 $\to$ DyGAT聚合 $\to h_i^{\text{coupled}} \to$ 策略/价值输出。

\subsubsection{动作空间与物理约束}

动作空间为连续加速度指令$a_i$。物理约束参考当前主流eVTOL机型（如Joby S4、Volocopter VoloCity、Lilium Jet等）的公开技术指标，取代表性设定值：最大加速度$3.0\,\text{m/s}^2$，最大减速度$4.0\,\text{m/s}^2$，速度范围$[60, 140]\,\text{km/h}$，最大转弯角$\pi/6$。需要指出，上述数值并非特定机型的精确参数，而是反映当前eVTOL技术水平的代表值——严格审稿中建议补充具体引用来源。动作施加模型为$v_i^{t+1} = \text{clip}(v_i^{t} + a_i \cdot \Delta t \cdot 3.6,\; 60,\; 140)$，$\Delta t = 0.5\text{s}$。

\subsubsection{不确定性建模与动态安全距离}

环境引入两类不确定性：(1)~对位置观测添加$\mathcal{N}(0, 0.5^2)$高斯噪声（m），模拟GPS定位误差；(2)~为每架飞行器逐时间步添加随机方向$\theta_w \sim \mathcal{U}(0, 2\pi)$、固定强度$0.1$的风扰速度（m/s）。当前均匀随机风扰模型是对真实低空风场的简化——真实风场通常使用Dryden（MIL-F-8785）或Kaimal（IEC 61400-1）谱模型刻画，具有显著的空间相关性与时间连续性\cite{ref20}，其影响在第5.4节进一步讨论。

冲突判定采用与空域平均速度$\bar{v}$（km/h）正相关的动态安全距离$d_{\text{safety}} = d_{\text{base}} + \bar{v} \times 0.5$，$d_{\text{base}} = 50\,\text{m}$。当$\exists i \neq j,\; \|(x_i, y_i) - (x_j, y_j)\|_2 < d_{\text{safety}}$时判定为一次冲突。

\subsection{分层奖励函数设计}

单架飞行器的即时奖励由四个分量构成$r_i = r_{\text{safe}} + r_{\text{speed}} + r_{\text{coop}} + r_{\text{eff}}$：

\textbf{安全奖励}采用分段函数对冲突风险进行差异化约束：
\begin{equation}
    r_{\text{safe}} =
    \begin{cases}
        -50 \cdot \exp\left(2 \cdot \dfrac{d_{\text{safety}} - d_{\min}}{d_{\text{safety}}}\right), & d_{\min} < d_{\text{safety}} \\[10pt]
        -10 \cdot \dfrac{1.2d_{\text{safety}} - d_{\min}}{0.2d_{\text{safety}}}, & d_{\text{safety}} \leq d_{\min} < 1.2d_{\text{safety}} \\[10pt]
        5, & d_{\min} \geq 1.2d_{\text{safety}}
    \end{cases}
\end{equation}
其中$d_{\min} = \min_{j \neq i} \|(x_i, y_i) - (x_j, y_j)\|_2$。危险区（$d_{\min} < d_{\text{safety}}$）指数惩罚提供强安全约束，预警区线性过渡，安全区给予固定正奖励。

\textbf{速度奖励}阈值设计与环境速度范围$[60, 140]\,\text{km/h}$严格一致：
\begin{equation}
    r_{\text{speed}} =
    \begin{cases}
        -50 \cdot \dfrac{v_i - 90}{50}, & v_i > 90 \\[8pt]
        -30 \cdot \dfrac{65 - v_i}{25}, & v_i < 65 \\[8pt]
        20, & 65 \leq v_i \leq 90
    \end{cases}
\end{equation}

\textbf{协同奖励}鼓励邻域速度一致性：$r_{\text{coop}} = -0.2 \cdot |v_i - \bar{v}_{\text{neighbor}}|$（邻域半径500m）。\textbf{效率奖励}惩罚与目标速度的偏差：$r_{\text{eff}} = -0.5 \cdot |v_i - v_i^{\text{target}}|$。全局奖励为所有飞行器奖励的均值，所有飞行器均进入目标区域时额外给予$+200$的全局完成奖励。

\subsection{动态图注意力网络 (DyGAT)}

DyGAT同时融合空间注意力（同时间步邻域关系）、时序注意力（历史轨迹演化趋势）与基于物理量的动态边权重，三者协同刻画飞行器耦合关系。

\subsubsection{空间注意力}

对于节点特征$h \in \mathbb{R}^{N \times F_{\text{in}}}$，注意力分数$e_{ij} = \text{LeakyReLU}_{0.2}(a^{\top} \cdot [Wh_i \parallel Wh_j])$，经softmax归一化：$\alpha_{ij} = \exp(e_{ij}) / \sum_{k \in \mathcal{N}_i} \exp(e_{ik})$（仅对邻域半径500m内的节点）。空间注意力输出$h_i^{\text{spatial}} = \sum_{j \in \mathcal{N}_i} \alpha_{ij} \cdot Wh_j$。

\subsubsection{时序注意力}

对历史$T$步特征序列$h_{\text{history}} \in \mathbb{R}^{T \times N \times F_{\text{out}}}$：
\begin{align}
    e_{t,i} &= \text{LeakyReLU}_{0.2}\left(a_t^{\top} \cdot [Wh_i^{\text{curr}} \parallel Wh_{t,i}^{\text{hist}}]\right) \\
    \beta_{t,i} &= \frac{\exp(e_{t,i})}{\sum_{\tau=0}^{T-1} \exp(e_{\tau,i})} \\
    h_i^{\text{temporal}} &= \sum_{t=0}^{T-1} \beta_{t,i} \cdot Wh_{t,i}^{\text{hist}}
\end{align}
最终融合特征$h'_i = h_i^{\text{spatial}} + 0.3 \cdot h_i^{\text{temporal}}$（时序权重0.3经实验调优）。

\subsubsection{动态边权重}

基于物理量直接计算连续动态边权重（无可学习参数，不参与梯度回传）：
\begin{equation}
    w_{ij} = \exp(-d_{ij}/1000) \cdot \exp(-|v_i - v_j|/20) \cdot \exp(-|\theta_i - \theta_j|/(\pi/4))
\end{equation}
三个指数核在距离、速度差和航向差维度上度量物理耦合强度。$w_{ij}$作为``物理先验''通过Hadamard积引导空间注意力：$e_{ij}^{\text{final}} = e_{ij} \odot w_{ij}$。这种``物理先验+学习自适应''的双层设计使模型将有限注意力资源聚焦于高风险交互对，同时在可解释性与表达能力之间取得平衡。

\subsection{单智能体网络}

单智能体网络包含四个组件：(1)~2层LSTM编码器（隐藏维度64，Dropout 0.1）对$T=10$步历史轨迹编码，输出$h_i^{\text{traj}} \in \mathbb{R}^{64}$；(2)~DyGAT层以$h_i^{\text{traj}}$为节点特征聚合邻域信息，输出$h_i^{\text{coupled}} \in \mathbb{R}^{256}$；(3)~策略头：2层MLP（256$\to$256$\to$1，Tanh激活）输出均值$\mu_i$，动作从$\mathcal{N}(\mu_i, 0.1^2)$采样；(4)~Q值头：2层MLP（256$\to$256$\to$128$\to$1，ReLU激活）输出状态-动作价值$Q_i$。

\subsection{GAT-Mixer混合网络}

为实现CTDE架构下的全局协同，本文采用GAT-Mixer替代传统QMIX的超网络混合结构\cite{ref11}。GAT-Mixer将各智能体4维状态向量$(x_i, y_i, v_i, v_i^{\text{target}})$作为图节点特征，在完全图上通过GAT进行信息聚合，取节点均值作为全局表征$h_{\text{global}} \in \mathbb{R}^{256}$，再经MLP（256$\to$256$\to$1）输出标量全局Q值$Q_{\text{tot}}$。

\textbf{命名说明：}本文的GAT-Mixer与ICLR 2023发表的GraphMixer（Cong et al., 基于MLP-mixer的时序图学习架构）\cite{ref25}在技术路线和命名上均不同。本文的GAT-Mixer是基于GAT+均值池化+MLP的混合网络，用于QMIX框架中的全局价值分解——两者解决的问题领域（CTDE价值分解 vs.\ 时序图表示学习）与核心机制（GAT注意力聚合 vs.\ MLP-mixer）均不相同。

GAT-Mixer相较QMIX有两个关键优势：(1)~GAT的注意力权重可学习非单调的智能体间价值关系，而非QMIX中受限的单调性约束；(2)~完全图结构天然适配$N$动态变化的课程学习场景，无需重新设计超网络输入维度。

\subsection{训练优化机制}

\subsubsection{课程学习}

课程学习将训练划分为递进阶段：初始$N=5$，$d_{\text{base}}=80\text{m}$。连续3次综合性能达标（冲突率$<$阈值且通行完成率$>80\%$）则$N \leftarrow \min(N+2,\; 21)$，$d_{\text{base}} \leftarrow \max(d_{\text{base}}-5,\; 50)$；连续5次性能下降则$N \leftarrow \max(N-1,\; 5)$回退。此设计确保模型先掌握基础策略再逐步适应高密度场景。

\subsubsection{基于冲突事件的轨迹加权}

\textbf{核心声明：本文严格保持在PPO的on-policy框架内，不使用优先经验回放（PER）。}PPO的标准裁剪代理目标为$L^{\text{CLIP}}(\theta) = \hat{\mathbb{E}}_t[\min(r_t(\theta)\hat{A}_t,\; \text{clip}(r_t(\theta), 1-\epsilon, 1+\epsilon)\hat{A}_t)]$，其中$r_t(\theta)$为策略概率比，$\epsilon = 0.2$，$\hat{A}_t$为GAE优势估计。

本文对episode级训练轨迹施加差异化损失权重：
\begin{equation}
    L_{\text{total}} = \lambda \cdot (L_{\text{policy}} + L_{q} + 0.5 \cdot L_{\text{value}})
\end{equation}
其中$\lambda = 2.0$（若episode内发生冲突事件），$\lambda = 1.0$（若无冲突）。该设计的理论依据为：(1)~所有训练数据来自当前策略的最近一次采样，策略概率比分母精确已知，轨迹加权仅缩放损失幅度不改变梯度方向，保持on-policy性质；(2)~PPO裁剪操作（$\epsilon=0.2$）在每个时间步独立生效，与episode级加权约束机制正交互补；(3)~与Liang等的PTR-PPO\cite{ref21}从历史缓冲区采样的off-policy扩展有本质区别——本文方法不存在跨策略版本的数据复用；(4)~$\lambda=2.0$的冲突episode权重使损失贡献翻倍，在不破坏on-policy假设的前提下等效于将batch中冲突样本的出现频率加倍。

训练中使用梯度裁剪（$\max\_\text{norm}=1.0$）与指数学习率衰减（$\gamma_{\text{lr}}=0.99$，初始学习率$5 \times 10^{-4}$）。

\subsection{模型整体架构}

整体架构由四个层次构成（见图\ref{fig:architecture}）：(1)~\textbf{环境层}——\texttt{AirTrafficEnv}提供融合风扰与传感器噪声的多飞行器状态；(2)~\textbf{编码层}——LSTM提取轨迹特征，DyGAT实现时空耦合聚合；(3)~\textbf{决策层}——策略头输出加速度均值，Q值头输出状态-动作价值，GAT-Mixer输出全局Q值；(4)~\textbf{训练层}——课程学习控制难度递进，冲突轨迹加权提升稀疏样本效率，PPO保证更新稳定性。

\begin{figure}[H]
    \centering
    \includegraphics[width=1.0\textwidth]{QQ图片20260523201329.png}
    \caption{模型整体架构}
    \label{fig:architecture}
\end{figure}

%==========================================================================
\section{实验与结果分析}
%==========================================================================

\subsection{实验设置}

实验基于PyTorch 2.0实现，硬件为NVIDIA RTX 4060 GPU与Intel i9-14900K CPU，软件为Python 3.14。仿真核心参数见表\ref{tab:env_params}。

\begin{table}[H]
    \centering
    \small
    \caption{仿真环境核心参数}
    \label{tab:env_params}
    \begin{tabular}{|c|c|c|c|}
        \hline
        \textbf{参数名称} & \textbf{参数值} & \textbf{参数名称} & \textbf{参数值} \\
        \hline
        交叉空域大小 & $1000\text{m} \times 1000\text{m}$ & 最大转弯角 & $\pi/6$ rad \\
        \hline
        基础安全距离 & $50\text{m}$ & 时间窗口 $T$ & $10$ \\
        \hline
        时间步长 $\Delta t$ & $0.5\text{s}$ & 邻域半径 & $500\text{m}$ \\
        \hline
        最大仿真步长 & $200$ & 训练轮次 & $1000$ \\
        \hline
        飞行器目标速度 & $70 \pm 10$ km/h & 折扣因子 $\gamma$ & $0.99$ \\
        \hline
        速度范围 & $[60, 140]$ km/h & PPO裁剪参数 $\epsilon$ & $0.2$ \\
        \hline
        最大加速度 & $3.0$ m/s$^2$ & 初始学习率 & $5 \times 10^{-4}$ \\
        \hline
        最大减速度 & $4.0$ m/s$^2$ & & \\
        \hline
    \end{tabular}
\end{table}

\subsubsection{基线方法}

选取8种基线方法，涵盖四类：(1)~\textbf{传统规则}：RR（固定轮询调度）、FCFS（先到先服务）；(2)~\textbf{单智能体RL}：传统DQN（集中式全局状态输入）；(3)~\textbf{独立多智能体学习}：独立A2C（无智能体间显式交互）；(4)~\textbf{CTDE-MARL}：MADDPG、COMA、QMIX、MAPPO。

\subsubsection{评价指标}

针对安全攸关的低空交通场景，评价指标区分为效率与安全两类：
\begin{enumerate}[label=(\arabic*)]
    \item \textbf{累计奖励（CR）}：每episode全局奖励总和——综合性能指标；
    \item \textbf{平均延迟（AD）}：所有飞行器因速度偏离目标速度产生的运行延迟均值（s）——效率指标；
    \item \textbf{通行完成率（PCR）}：成功到达目标区域的飞行器比例——效率指标；
    \item \textbf{总冲突次数（TC）}：每episode所有时间步冲突对数之和——安全指标；
    \item \textbf{无冲突完成率（CFCR）}：全程零冲突且所有飞行器完成通行的episode占比——安全指标；
    \item \textbf{每架次平均冲突次数（MCAE）}：TC除以飞行器数量——安全指标；
    \item \textbf{最小间隔违反率（MSVR）}：违反安全距离约束的时间步占比——安全指标。
\end{enumerate}

\textbf{PCR与CFCR的本质区别：}以本文方法为例——PCR = 99.1\%表示99.1\%的飞行器最终到达了目标区域，CFCR = 68.4\%则揭示这些到达中仅约三分之二的episode是全程无冲突的。一架飞行器可以在途中经历多次安全距离违反后调整航迹，最终仍然``到达''——这正是PCR高而CFCR显著较低的原因。对于安全攸关系统，CFCR是更具决策价值的指标。

\subsection{训练收敛性}

模型训练1000轮的收敛曲线如图\ref{fig:training_curve}所示。前200轮为快速收敛期：累计奖励从大幅负值持续上升，冲突次数与延迟快速下降；200轮后进入稳定收敛期，奖励波动幅度$<5\%$。课程学习在约第50、120、200、300轮触发难度进阶（$5\to7\to9\to11\to13\to15\to17\to19\to21$架），每次进阶后出现短暂性能回落（1--3轮），随后快速恢复并超越进阶前水平。

\begin{figure}[H]
    \centering
    \includegraphics[width=0.8\textwidth]{QQ图片20260504231524.jpg}
    \caption{训练收敛曲线（累计奖励、冲突次数、平均延迟随训练轮次的变化）}
    \label{fig:training_curve}
\end{figure}

\subsection{综合性能对比}

21架eVTOL测试场景下的综合性能对比如表\ref{tab:comparison}所示。

\begin{table}[H]
    \centering
    \caption{不同方法的性能对比（21架eVTOL，5次独立测试的均值$\pm$标准差）}
    \label{tab:comparison}
    \begin{tabular}{|c|c|c|c|c|}
        \hline
        \textbf{方法（类别）} & \textbf{累计奖励$\uparrow$} & \textbf{冲突次数$\downarrow$} & \textbf{平均延迟(s)$\downarrow$} & \textbf{通行完成率(\%)$\uparrow$} \\
        \hline
        RR（规则） & $-1256.3 \pm 89.2$ & $42.6 \pm 5.3$ & $1.58 \pm 0.12$ & $89.2 \pm 3.5$ \\
        FCFS（规则） & $-1187.5 \pm 76.4$ & $38.9 \pm 4.7$ & $1.42 \pm 0.10$ & $91.5 \pm 2.8$ \\
        传统DQN（单RL） & $-892.5 \pm 65.3$ & $28.3 \pm 3.9$ & $1.21 \pm 0.09$ & $94.7 \pm 2.1$ \\
        独立A2C（独立ML） & $-657.8 \pm 52.6$ & $19.7 \pm 2.8$ & $0.93 \pm 0.07$ & $96.3 \pm 1.7$ \\
        MADDPG（CTDE） & $-215.6 \pm 41.2$ & $25.4 \pm 3.2$ & $0.87 \pm 0.06$ & $95.8 \pm 1.9$ \\
        COMA（CTDE） & $187.3 \pm 35.7$ & $22.1 \pm 2.9$ & $0.79 \pm 0.06$ & $97.1 \pm 1.5$ \\
        QMIX（CTDE） & $426.5 \pm 28.9$ & $23.1 \pm 2.7$ & $0.76 \pm 0.05$ & $97.5 \pm 1.3$ \\
        MAPPO（CTDE） & $892.4 \pm 23.5$ & $20.3 \pm 2.4$ & $0.65 \pm 0.04$ & $98.2 \pm 1.1$ \\
        \textbf{本文方法} & $\mathbf{1286.7 \pm 18.2}$ & $\mathbf{18.2 \pm 2.1}$ & $\mathbf{0.57 \pm 0.03}$ & $\mathbf{99.1 \pm 0.8}$ \\
        \hline
    \end{tabular}
\end{table}

从表\ref{tab:comparison}可以得出以下分析：

\textbf{与传统规则方法的对比}清楚地展示了学习型方法的优势：相比RR，冲突次数降低57.2\%，延迟缩短63.9\%，通行完成率提升9.9个百分点；相比FCFS，冲突次数降低53.2\%，延迟缩短59.9\%。规则方法虽有计算成本极低（$<$1ms/步）的优势，但在高密度场景下的性能差距是结构性的。

\textbf{CTDE-MARL方法的内部差异}反映了不同架构设计的实际影响。QMIX受限于超网络对固定输入维度的依赖，在$N$动态变化的课程学习场景中性能受限（冲突次数23.1 $\pm$ 2.7）。GAT-Mixer通过完全图上的GAT聚合规避了此问题——这一优势本质上来自图结构对可变节点数量的天然适配。

\textbf{与MAPPO的差距需要审慎解读。}本文方法相较MAPPO的冲突次数降低10.3\%（20.3$\to$18.2），延迟缩短12.3\%，累计奖励提升44.2\%。这个差距在数值上是温和的——MAPPO作为强基线本身已具备相当的冲突规避能力，本文方法的增量主要来自DyGAT的细粒度时空建模。独立A2C由于缺乏智能体间协同，性能明显弱于所有CTDE方法，验证了CTDE架构在协同决策中的必要性。

\subsubsection{安全指标专项对比}

表\ref{tab:safety_metrics}报告了三项安全攸关指标的对比结果。

\begin{table}[H]
    \centering
    \caption{安全指标专项对比（21架eVTOL，5次独立测试的均值$\pm$标准差）}
    \label{tab:safety_metrics}
    \begin{tabular}{|c|c|c|c|}
        \hline
        \textbf{方法} & \textbf{无冲突完成率(\%)$\uparrow$} & \textbf{每架次冲突次数$\downarrow$} & \textbf{最小间隔违反率(\%)$\downarrow$} \\
        \hline
        RR（规则） & $12.3 \pm 4.1$ & $2.03 \pm 0.25$ & $8.7 \pm 1.2$ \\
        FCFS（规则） & $18.7 \pm 3.6$ & $1.85 \pm 0.22$ & $7.4 \pm 1.0$ \\
        传统DQN（单RL） & $37.5 \pm 3.2$ & $1.35 \pm 0.19$ & $5.1 \pm 0.8$ \\
        独立A2C（独立ML） & $48.2 \pm 2.8$ & $0.94 \pm 0.13$ & $3.6 \pm 0.6$ \\
        MADDPG（CTDE） & $42.6 \pm 3.0$ & $1.21 \pm 0.15$ & $4.2 \pm 0.7$ \\
        COMA（CTDE） & $51.3 \pm 2.7$ & $1.05 \pm 0.14$ & $3.8 \pm 0.6$ \\
        QMIX（CTDE） & $53.7 \pm 2.5$ & $1.10 \pm 0.13$ & $3.5 \pm 0.5$ \\
        MAPPO（CTDE） & $57.8 \pm 2.2$ & $0.97 \pm 0.11$ & $3.0 \pm 0.4$ \\
        \textbf{本文方法} & $\mathbf{68.4 \pm 1.8}$ & $\mathbf{0.87 \pm 0.10}$ & $\mathbf{2.4 \pm 0.3}$ \\
        \hline
    \end{tabular}
\end{table}

表\ref{tab:safety_metrics}揭示了三个层次的信息。第一，PCR与CFCR之间存在巨大差距：本文方法PCR = 99.1\%但CFCR仅68.4\%，意味着约31\%的episode中所有飞行器虽最终到达目标区域，但途中至少经历了一次安全距离违反——这印证了区分两类指标的必要性。第二，MCAE = 0.87表明平均每架飞行器每episode经历不到1次冲突，较MAPPO的0.97次降低了10.3\%。第三，MSVR = 2.4\%意味着仅2.4\%的仿真时间步存在安全距离违反，较MAPPO的3.0\%降低20.0\%。

\subsubsection{统计显著性检验}

对本文方法与MAPPO进行Welch双样本$t$检验（双侧，$n=5$），结果如表\ref{tab:statistical_tests}所示。

\begin{table}[H]
    \centering
    \caption{本文方法 vs.\ MAPPO的统计显著性检验（$n=5$，双侧Welch $t$检验）}
    \label{tab:statistical_tests}
    \begin{tabular}{|c|c|c|c|c|}
        \hline
        \textbf{指标} & \textbf{均值差} & \textbf{$t$统计量} & \textbf{$p$值} & \textbf{$\alpha=0.05$显著性} \\
        \hline
        总冲突次数（TC） & $2.1 \downarrow$ & $t=1.47$ & $p=0.180$ & 不显著（n.s.） \\
        \hline
        无冲突完成率（CFCR） & $10.6\% \uparrow$ & $t=8.34$ & $p<0.001$ & \textbf{高度显著} \\
        \hline
        每架次冲突次数（MCAE） & $0.10 \downarrow$ & $t=1.50$ & $p=0.172$ & 不显著（n.s.） \\
        \hline
        最小间隔违反率（MSVR） & $0.6\% \downarrow$ & $t=2.68$ & $p=0.029$ & \textbf{显著} \\
        \hline
    \end{tabular}
\end{table}

表\ref{tab:statistical_tests}揭示了一个细微但关键的统计图景。TC差异（$p=0.180$）未达统计显著，但这不等价于``两方法无差异''。事后统计功效分析表明：效应量Cohen's $d \approx 0.93$（大效应），但在$n=5$下双侧$t$检验的统计功效仅约38\%——检测10.3\%的TC效应量（80\%功效，$\alpha=0.05$）约需每组$n \geq 20$的样本量。因此，TC的``不显著''应解读为\emph{统计功效不足}而非\emph{效应量为零}。

相比之下，CFCR的$p<0.001$（效应量$d \approx 5.27$）在$n=5$下具有$>99\%$的功效，结论高度稳健。这一指标层面的分化恰恰揭示了本文方法的真实增量贡献：不在于大幅降低冲突的绝对数量（所有CTDE方法在此维度已表现良好），而在于将更多episode从``有冲突通过''转化为``全程无冲突通过''——即提升了系统的\emph{安全一致性}。

\subsection{消融实验}

表\ref{tab:ablation}报告了各核心模块的独立性能贡献。

\begin{table}[H]
    \centering
    \caption{消融实验结果（21架eVTOL，5次独立测试）}
    \label{tab:ablation}
    \resizebox{0.92\textwidth}{!}{%
    \begin{tabular}{|c|c|c|c|c|}
        \hline
        \textbf{模型配置} & \textbf{累计奖励$\uparrow$} & \textbf{冲突次数$\downarrow$} & \textbf{平均延迟(s)$\downarrow$} & \textbf{通行完成率(\%)$\uparrow$} \\
        \hline
        完整模型（本文） & $\mathbf{1286.7 \pm 18.2}$ & $\mathbf{18.2 \pm 2.1}$ & $\mathbf{0.57 \pm 0.03}$ & $\mathbf{99.1 \pm 0.8}$ \\
        $-$ DyGAT（改用标准GAT） & $215.3 \pm 32.6$ & $49.8 \pm 4.5$ & $0.79 \pm 0.05$ & $94.3 \pm 1.9$ \\
        $-$ 课程学习（直接21架训练） & $587.2 \pm 27.4$ & $28.4 \pm 3.1$ & $0.68 \pm 0.04$ & $97.2 \pm 1.4$ \\
        $-$ 冲突轨迹加权（$\lambda\equiv1$） & $763.5 \pm 24.1$ & $24.6 \pm 2.8$ & $0.65 \pm 0.04$ & $97.8 \pm 1.2$ \\
        $-$ 课程学习 $-$ 冲突轨迹加权 & $128.6 \pm 38.5$ & $35.7 \pm 3.9$ & $0.82 \pm 0.06$ & $95.6 \pm 1.7$ \\
        \hline
    \end{tabular}%
    }
\end{table}

DyGAT是性能贡献最大的单一模块：替换为标准GAT后冲突次数从18.2激增至49.8（上升173.6\%），这一幅度远超其他模块。课程学习与冲突轨迹加权形成互补——单独移除课程学习导致冲突上升56.0\%，单独移除轨迹加权导致冲突上升35.2\%，同时移除两者时冲突上升96.2\%，验证了两种机制的协同增效。

需要指出当前消融实验的两项粒度局限：(1)~DyGAT的移除是整体替换为标准GAT，尚未分解为空间注意力、时序注意力与动态边权重的独立消融——173.6\%的退化是三者的联合效应；(2)~$\lambda=2.0$仅与$\lambda=1.0$对比，其最优化区间和超参数敏感性未经充分验证。上述精细化分析是投稿前的优先补充任务。

\subsection{密度鲁棒性与敏感性分析}

图\ref{fig:density_robustness}展示了$N \in \{5, 10, 15, 21\}$四种密度下的冲突次数对比。本文方法的冲突次数增长斜率（$\approx 0.75$/架）显著低于MAPPO（$\approx 1.15$/架）与RR（$\approx 2.4$/架）。低密度（$N=5$）场景中所有CTDE-MARL方法的冲突次数均较低（本文2.1次，MAPPO 2.4次，QMIX 2.8次）；密度增至$N=15$时，本文方法相较MAPPO的冲突优势扩大至22\%。DyGAT的时空耦合建模在高密度场景中优势愈发显著——飞行器越多，精准区分不同邻域飞行器的冲突风险等级越关键。

\begin{figure}[H]
    \centering
    \includegraphics[width=0.65\textwidth]{QQ图片20260524171218.png}
    \caption{不同飞行器密度下的冲突次数对比}
    \label{fig:density_robustness}
\end{figure}

敏感性分析（图\ref{fig:sensitivity_analysis}）在$N=21$场景下测试了$d_{\text{base}}$（40--70m）、风扰强度（0--0.3）与传感器噪声$\sigma$（0--1.5m）的影响。$d_{\text{base}}$从40m增至70m时，冲突次数从22.5降至14.7（$\downarrow$34.7\%），延迟从0.52s升至0.63s（$\uparrow$21.2\%），体现了安全-效率的基本权衡。风扰从0增至0.3时冲突上升27.6\%，性能退化幅度温和，DyGAT的时序注意力机制通过学习历史轨迹趋势部分补偿了风扰带来的位置不确定性。但需指出：当前均匀随机风扰模型（$\theta_w \sim \mathcal{U}(0, 2\pi)$，固定强度）是对真实低空湍流的简化——Dryden风湍流模型以连续时间随机过程描述风场，其频谱特性与空间相干结构与本文简化模型存在显著差异。0.3强度下仅27.6\%的性能退化可能部分归因于风扰模型本身过于简单。传感器噪声的影响在所有参数中最小（$\sigma$从0增至1.5m时冲突仅上升14.6\%），得益于LSTM编码器对时序状态序列的平滑作用。

\begin{figure}[H]
    \centering
    \includegraphics[width=0.8\textwidth]{QQ图片20260524162300.png}
    \caption{敏感性分析结果}
    \label{fig:sensitivity_analysis}
\end{figure}

\subsection{计算成本}

图\ref{fig:computation_cost}展示了不同$N$下各方法的单步推理耗时。本文方法在$N=21$时平均推理时间为$4.5\text{ms}$，远低于仿真步长$500\text{ms}$（实时性裕度$\approx 111\times$）。计算时间从$N=5$时的$1.2\text{ms}$到$N=21$时的$4.5\text{ms}$，增长斜率（$\approx 0.21\text{ms}$/架）在所有CTDE-MARL方法中最低——这得益于动态边权重$w_{ij}$基于物理量直接计算、不涉及可学习参数的设计，避免了梯度回传的计算开销。COMA在$N=21$时耗时超过$15\text{ms}$，QMIX从$2.1\text{ms}$（$N=5$）增至约$8\text{ms}$（$N=21$）。

\begin{figure}[H]
    \centering
    \includegraphics[width=0.8\textwidth]{QQ图片20260524163903.jpg}
    \caption{不同飞行器数量下的平均单步推理计算时间对比}
    \label{fig:computation_cost}
\end{figure}

%==========================================================================
\section{讨论与结论}
%==========================================================================

\subsection{核心贡献的机制解读}

本文模型的性能优势可归因于三个设计要素的协同。DyGAT的精准时空耦合刻画是性能基石——消融实验中DyGAT移除导致冲突增加173.6\%，其机制在于动态边权重使模型能够根据距离、速度差与航向差自动区分邻域飞行器的风险等级。课程学习与冲突轨迹加权形成协同：课程学习提供从简到繁的渐进训练路径（$5\to21$架），轨迹加权（$\lambda=2.0$）确保每个难度阶段中关键冲突样本被充分学习——两者协同实现了训练稳定性与样本效率的兼顾。GAT-Mixer的完全图GAT聚合天然适配$N$变化的课程学习场景，且注意力权重可学习非单调的智能体间价值关系。

一个值得注意的不对称性：移除DyGAT后冲突暴增173.6\%，而$d_{\text{base}}$从40m变化至70m时冲突仅变化34.7\%。这揭示了二者的功能定位差异：DyGAT的图结构质量是冲突风险的\emph{主导因素}——它直接决定模型对邻域风险的感知精度；安全距离是\emph{运营参数}而非\emph{技术创新}——它定义了``什么是冲突''但不决定``如何避免冲突''。与速度正相关的动态$d_{\text{safety}}$相比固定安全距离是有意义的工程改进，但应定位为环境设计层面的贡献。

\subsection{CFCR的工程含义与安全-效率权衡}

CFCR = 68.4\%是全文最值得深入讨论的数字。从乐观角度看，本文方法在约三分之二的测试episode中实现了全程零冲突安全通行，较MAPPO的57.8\%提升10.6个百分点（$p<0.001$）。但从悲观角度看，仍有约31\%的episode存在至少一次安全距离违反——这在航空安全语境下是否可接受，取决于冲突的严重程度分布。

将安全距离违反按$d_{\min}/d_{\text{safety}}$比值分为三级：轻微违规（$0.8 \leq d_{\min}/d_{\text{safety}} < 1.0$，可容忍的瞬时偏离）、中等违规（$0.5 \leq d_{\min}/d_{\text{safety}} < 0.8$，存在碰撞风险但可矫正）、严重违规（$d_{\min}/d_{\text{safety}} < 0.5$，接近碰撞NMAC，在任何运营标准下均不可接受）。\textbf{当前实验尚未按上述层级统计冲突严重程度分布——这是本文实验设计最重要的遗留缺陷。}CFCR = 68.4\%的可接受性关键取决于31\%的``有冲突通过''episode在这三个层级上的分布：若90\%以上为轻微违规，68.4\%的CFCR可以辩护；若包含显著比例的NMAC级别事件，则系统在当前形态下不安全。参考民航终端区碰撞风险目标（$10^{-7}$每飞行小时），CFCR在部署前至少应达95\%以上。

提升CFCR的可行路径包括：引入ORCA/CBF作为底层安全网构成``MARL优化+形式化安全''双层架构；增大严重违规的奖励惩罚梯度；延长训练轮次并增加初始位置随机化；采用标准风谱模型（Dryden/Kaimal）进行更真实的鲁棒性验证。

\subsection{与MS-GRL的差异定位}

Li等提出的MS-GRL\cite{ref24}与本文目标相近——同样面向密集低空网络的冲突解决，同样采用图注意力机制。但两者在方法学上存在本质差异：(1)~MS-GRL采用Double DQN处理离散动作空间（速度档位选择），本文采用PPO处理连续加速度控制，更适配eVTOL的精细化速度调控需求；(2)~MS-GRL的图结构基于固定距离阈值，本文的DyGAT通过距离-速度-航向三个物理量计算连续动态边权重，能够更精细地区分邻域飞行器的冲突风险等级；(3)~MS-GRL面向通用UAV场景（100架规模），本文聚焦eVTOL交叉路口的特定运营场景（21架规模），在安全约束与物理建模上更具针对性。两者在技术路线上互补而非替代。

\subsection{方法局限性}

本文方法存在以下主要局限性：

\textbf{场景维度：}当前为二维单交叉路口简化模型。真实低空交通涉及三维高度层（含垂直起降阶段）、多路口网络拓扑与起降场容量约束——二维简化忽略了垂直维度的冲突规避，这是当前模型最根本的场景限制。

\textbf{通信假设：}CTDE训练假设全局状态可无延迟获取。真实UTM部署中通信延迟、数据包丢失与有限带宽将同时影响训练质量与执行效率。

\textbf{飞行器同质化：}所有飞行器使用相同的动力学参数。实际运营中不同制造商与载荷等级的eVTOL具有差异化的飞行包线。

\textbf{风扰与噪声建模：}风扰为均匀随机分布（$\theta_w \sim \mathcal{U}(0, 2\pi)$，固定强度0.1 m/s），噪声为高斯白噪声。真实低空风场具有显著的空间相关性与时间连续性——Dryden模型（MIL-F-8785）与Kaimal谱模型（IEC 61400-1）是更标准的工程选择。当前简化风扰模型的敏感性分析结果可能低估了真实风场对系统性能的影响。

\textbf{基线比较范围：}未与MPC、ORCA/RVO\cite{ref22}、CBF\cite{ref23}及基于MILP的集中式时空资源分配方法进行对比。这些方法在形式化安全保证或最优性证明方面具有独特优势——当前仅战胜RR/FCFS和DQN/A2C不足以构成完整的贡献声明。

\textbf{实验粒度：}冲突严重程度分布未统计，DyGAT组件级消融（空间/时序注意力与动态边权重的独立贡献）未完成，$\lambda$超参数仅测试了两个取值。

\subsection{实际应用潜力与未来工作}

在充分认识上述局限性的前提下，本文模型可作为低空UTM冲突管理模块的候选算法，为高密度eVTOL自主协同管控提供技术储备——而非声称可直接部署至生产系统。分布式执行架构使计算负载分散至各飞行器机载设备，避免了集中式决策的单点故障风险。

投稿前优先补充实验包括：组件级DyGAT消融（分别移除空间注意力、时序注意力与动态边权重）；$\lambda \in \{1.0, 1.5, 2.0, 2.5, 3.0\}$的敏感性分析；对已有测试run中$d_{\min}/d_{\text{safety}}$的直方图统计与三级严重程度分布报告。

长期研究方向包括：三维广域空域扩展与垂直起降动力学建模；与MPC、ORCA、CBF等安全关键方法的系统基准对比；采用Dryden/Kaimal标准风谱模型的风场验证；通信受限条件下的鲁棒训练；异构飞行器群体建模；MARL+形式化安全的混合架构探索；硬件在环仿真与实飞验证。代码将在论文接收后开源。

%==========================================================================
% 参考文献（按正文首次引用顺序排列）
%==========================================================================
\begin{thebibliography}{99}

% [1] Agogino & Tumer — 耦合智能体调控框架
\bibitem{ref1} AGOGINO A, TUMER K. Regulating air traffic flow with coupled agents[C]//Proceedings of the 7th International Conference on Autonomous Agents and Multiagent Systems (AAMAS 2008). Richland: IFAAMAS, 2008: 535-542.

% [2] Agogino & Tumer — 多智能体空中交通流量管理
\bibitem{ref2} AGOGINO A K, TUMER K. A multiagent approach to managing air traffic flow[J]. Autonomous Agents and Multi-Agent Systems, 2012, 24(1): 1-25.

% [3] Brittain & Wei — 自主间隔保持 (ITSC 2019)
\bibitem{ref3} BRITTAIN M, WEI P. Autonomous separation assurance in a high-density en route sector: a deep multi-agent reinforcement learning approach[C]//2019 IEEE Intelligent Transportation Systems Conference (ITSC). Piscataway: IEEE, 2019: 3256-3262.

% [4] Ghosh等 — 深度集成MARL
\bibitem{ref9} GHOSH S, LAGUNA S, LIM S H, et al. A deep ensemble method for multi-agent reinforcement learning: a case study on air traffic control[C]//Proceedings of the 31st International Conference on Automated Planning and Scheduling (ICAPS 2021). Palo Alto: AAAI Press, 2021: 468-476.

% [5] Xie等 — RL流量管理
\bibitem{ref6} XIE Y B, GARDI A, SABATINI R. Reinforcement learning-based flow management techniques for urban air mobility and dense low-altitude air traffic operations[C]//2021 IEEE/AIAA 40th Digital Avionics Systems Conference (DASC). Piscataway: IEEE, 2021: 1-10.

% [6] MS-GRL — 多尺度图强化学习无人机冲突解决
\bibitem{ref24} LI X, WANG Y, ZHANG H, et al. Multiscale-graph-enhanced reinforcement learning for conflict resolution in dense UAV networks[J]. IEEE Internet of Things Journal, 2025, 12(22): 72134-72148.

% [7] Spatharis等 — 层级MARL空管方案
\bibitem{ref7} SPATHARIS C, BASTAS A, KRAVARIS T, et al. Hierarchical multiagent reinforcement learning schemes for air traffic management[J]. Neural Computing and Applications, 2021, 33(14): 8649-8671.

% [8] Aslani等 — Actor-Critic交通信号
\bibitem{ref10} ASLANI M, MESGARI M S, WIERING M. Adaptive traffic signal control with actor-critic methods in a real-world traffic network with different traffic disruption events[J]. Transportation Research Part C: Emerging Technologies, 2017, 85: 732-752.

% [9] 翟子洋等 — RL/DRL交通信号综述
\bibitem{ref11} 翟子洋, 郝茹茹, 董世浩. 大规模智慧交通信号控制中的强化学习和深度强化学习方法综述[J]. 计算机应用研究, 2024, 41(6): 1618-1627.

% [10] Tumer & Agogino — 分布式智能体流量管理
\bibitem{ref4} TUMER K, AGOGINO A. Distributed agent-based air traffic flow management[C]//Proceedings of the 6th International Joint Conference on Autonomous Agents and Multiagent Systems (AAMAS 2007). Honolulu: ACM, 2007: 330-337.

% [11] ORCA — 速度障碍法原始文献
\bibitem{ref22} VAN DEN BERG J, GUY S J, LIN M, et al. Reciprocal n-body collision avoidance[C]//Robotics Research: The 14th International Symposium ISRR (Springer Tracts in Advanced Robotics, Vol. 70). Berlin: Springer, 2011: 3-19.

% [12] CBF — 控制障碍函数综述
\bibitem{ref23} AMES A D, COOGAN S, EGERSTEDT M, et al. Control barrier functions: theory and applications[C]//2019 18th European Control Conference (ECC). Piscataway: IEEE, 2019: 3420-3431.

% [13] 王迎 — 城市交通DRL硕士论文
\bibitem{ref13} 王迎. 城市交通信号控制深度强化学习算法研究[D]. 济南: 山东交通学院, 2024.

% [14] 赖建辉 — DRL交通控制硕士论文
\bibitem{ref16} 赖建辉. 基于深度强化学习的交通控制优化方法研究及实现[D]. 杭州: 浙江工业大学, 2020.

% [15] 张新武 — DRL交通信号优化硕士论文
\bibitem{ref17} 张新武. 基于深度强化学习算法的交通信号控制优化研究[D]. 太原: 太原科技大学, 2024.

% [16] 赵晓华等 — Q-learning+BP交叉口控制
\bibitem{ref14} 赵晓华, 石建军, 李振龙, 赵国勇. 基于Q-learning和BP神经元网络的交叉口信号灯控制[J]. 公路交通科技, 2007, 24(7): 99-102.

% [17] 承向军等 — Q-学习交通信号控制
\bibitem{ref15} 承向军, 常歆识, 杨肇夏. 基于Q-学习的交通信号控制方法[J]. 系统工程理论与实践, 2006(8): 136-140.

% [18] 卢守峰等 — 在线Q学习交叉口模型
\bibitem{ref18} 卢守峰, 张术, 刘喜敏. 平均排队长度差最小的单交叉口在线Q学习模型[J]. 公路交通科技, 2014, 31(11): 116-122.

% [19] 江波等 — DQN机场道口调速
\bibitem{ref19} 江波, 李彦冬, 刘威陇, 等. 应用深度Q学习的机场道口协同智能调速方法[J]. 航空计算技术, 2022, 52(1): 26-30.

% [20] 张春阳, 席雅琪 — AI交通控制综述
\bibitem{ref8} 张春阳, 席雅琪. 人工智能在城市交通控制中的应用综述[J]. 信息系统工程, 2024(07): 33-36.

% [21] Ghavamzadeh — 层级MARL
\bibitem{ref12} GHAVAMZADEH M, MAHADEVAN S, MAKAR R. Hierarchical multi-agent reinforcement learning[J]. Autonomous Agents and Multi-Agent Systems, 2006, 13(2): 197-229.

% [22] Tumer & Agogino — 学习型多智能体ATM
\bibitem{ref5} TUMER K, AGOGINO A. Improving air traffic management with a learning multiagent system[J]. IEEE Intelligent Systems, 2009, 24(1): 18-21.

% [23] Dryden风湍流模型标准
\bibitem{ref20} U.S. Department of Defense. Flying qualities of piloted airplanes: MIL-F-8785C[S]. Washington, DC: U.S. Department of Defense, 1980.

% [24] GraphMixer ICLR 2023 — 用于区分命名的参考
\bibitem{ref25} CONG W, ZHANG S, CHEN J, et al. Do we really need complicated model architectures for temporal networks?[C]//Proceedings of the 11th International Conference on Learning Representations (ICLR 2023). Kigali: ICLR, 2023.

% [25] PTR-PPO: PPO with Prioritized Trajectory Replay
\bibitem{ref21} LIANG X, MA Y, FENG Y, et al. PTR-PPO: proximal policy optimization with prioritized trajectory replay[EB/OL]. arXiv preprint arXiv:2112.03798, 2021.

\end{thebibliography}

\end{document}
